% Bagian: first-order-logic
% Bab: natural-deduction
% Bagian: quantifier-rules

\documentclass[../../../include/open-logic-section]{subfiles}

\begin{document}

\olfileid[id]{fol}{ntd}{qrl}

\olsection{Aturan Kuantor}

\subsection{Aturan untuk $\lforall$}

\begin{defish}
\AxiomC{$!A(a)$}
\RightLabel{\Intro{\lforall}}
\UnaryInfC{$\lforall[x][\Atom{!A}{x}]$}
\DisplayProof
\hfill
\AxiomC{$\lforall[x][\Atom{!A}{x}]$}
\RightLabel{\Elim{\lforall}}
\UnaryInfC{$!A(t)$}
\DisplayProof
\end{defish}

Dalam aturan untuk~$\lforall$, $t$ adalah suku tertutup (suku yang tidak
memuat variabel), dan $a$~adalah !!a{constant} yang tidak muncul dalam
kesimpulan~$\lforall[x][!A(x)]$, ataupun dalam asumsi mana pun yang
!!{undischarged} dalam !!{derivation} yang berakhir dengan premis~$!A(a)$.
Kita menyebut $a$ sebagai \emph{variabel eigen} dari inferensi
\Intro{\lforall}.\footnote{Kita menggunakan istilah ``variabel eigen''
meskipun $a$ dalam aturan di atas adalah !!a{constant}. Alasannya
bersifat historis.}

\subsection{Aturan untuk $\lexists$}

\begin{defish}
\AxiomC{$\Atom{!A}{t}$}
\RightLabel{\Intro{\lexists}}
\UnaryInfC{$\lexists[x][\Atom{!A}{x}]$}
\DisplayProof
\hfill
\AxiomC{$\lexists[x][\Atom{!A}{x}]$}
\AxiomC{[$\Atom{!A}{a}$]$^n$}
\DeduceC{$!C$}
\DischargeRule{\Elim{\lexists}}{n}
\BinaryInfC{$!C$}
\DisplayProof
\end{defish}

Sekali lagi, $t$ adalah suku tertutup, dan $a$ adalah !!a{constant}
yang tidak muncul dalam premis $\lexists[x][!A(x)]$, dalam
kesimpulan~$!C$, ataupun dalam asumsi mana pun yang !!{undischarged}
dalam !!{derivation}s yang berakhir dengan kedua premis tersebut
(selain asumsi $!A(a)$). Kita menyebut $a$ sebagai \emph{variabel
eigen} dari inferensi \Elim{\lexists}.

Pembatasan atas kemunculan variabel eigen dalam inferensi
\Intro{\lforall} atau \Elim{\lexists} yang baru saja dinyatakan disebut
\emph{syarat variabel eigen}.

\begin{explain}
Ingat kembali konvensi bahwa ketika $!A$ adalah !!a{formula} dengan
!!{variable}~$x$ bebas, kita menandainya dengan menulis~$!A(x)$. Dalam
konteks yang sama, $!A(t)$ adalah singkatan dari~$\Subst{!A}{t}{x}$.
Jadi, aturan $\Intro\lexists$ juga dapat kita tulis sebagai:
\begin{prooftree}
\AxiomC{$\Subst{!A}{t}{x}$}
\RightLabel{\Intro{\lexists}}
\UnaryInfC{$\lexists[x][!A]$}
\end{prooftree}
Perhatikan bahwa $t$ mungkin sudah muncul dalam~$!A$; misalnya, $!A$
mungkin berupa~$\Atom{\Obj P}{t,x}$. Jadi, menyimpulkan
$\lexists[x][\Atom{\Obj P}{t,x}]$ dari~$\Atom{\Obj P}{t,t}$ merupakan
penerapan $\Intro\lexists$ yang benar---kita boleh ``mengganti'' satu
atau lebih, dan tidak harus semua, kemunculan~$t$ dalam premis dengan
!!{variable}~$x$ yang terikat. Akan tetapi, syarat variabel eigen dalam
$\Intro\lforall$ dan~$\Elim\lexists$ mengharuskan !!{constant}~$a$
tidak muncul dalam~$!A$. Karena itu, kita tidak dapat menyimpulkan
$\lforall[x][\Atom{\Obj P}{a,x}]$ secara benar dari
$\Atom{\Obj P}{a,a}$ dengan menggunakan~$\Intro\lforall$.
\end{explain}

\begin{explain}
Dalam \Intro{\lexists} dan \Elim{\lforall} tidak ada pembatasan variabel
eigen, dan suku~$t$ dapat berupa sebarang suku tertutup, sehingga tidak
ada syarat semacam itu yang perlu dikhawatirkan. Sebaliknya, dalam aturan
\Elim{\lexists} dan \Intro{\lforall}, syarat variabel eigen mengharuskan
!!{constant}~$a$ tidak muncul di mana pun dalam kesimpulan atau dalam
asumsi yang !!{undischarged}. Syarat ini diperlukan untuk memastikan
bahwa sistem tersebut sahih, yakni hanya !!{derive} !!{sentence}s dari
asumsi-asumsi !!{undischarged} yang mengakibatkannya. Tanpa syarat ini,
hal berikut akan diizinkan:
\begin{prooftree}
  \AxiomC{$\lexists[x][!A(x)]$}
  \AxiomC{$\Discharge{!A(a)}{1}$}
  \RightLabel{*\Intro{\lforall}}
  \UnaryInfC{$\lforall[x][!A(x)]$}
  \RightLabel{\Elim{\lexists}}
  \BinaryInfC{$\lforall[x][!A(x)]$}
\end{prooftree}
Padahal, $\lexists[x][!A(x)] \Entails/ \lforall[x][!A(x)]$.

Karena aturan-aturan eliminasi kuantor hanya mengizinkan suku tertutup untuk
menggantikan !!{variable}s, setiap !!{formula} yang dapat diturunkan
dari suatu himpunan !!{sentence}s dengan sendirinya merupakan
!!a{sentence}.
\end{explain}


\end{document}
